\documentclass[11pt]{article}
\usepackage[a4paper,margin=27mm]{geometry}
\usepackage[T1]{fontenc}
\usepackage{lmodern,microtype}
\usepackage{amsmath,amssymb,amsthm,mathtools}
\usepackage{booktabs,array}
\usepackage[hidelinks]{hyperref}
\hypersetup{pdftitle={Sharp worst-case factors for randomly pivoted Cholesky and LU},
 pdfsubject={Proposed negative resolutions of the same-rank RPCholesky and RPLU conjectures},
 pdfkeywords={randomized numerical linear algebra, Cholesky, LU, sharp bounds}}
\usepackage{enumitem}
\setlength{\parskip}{4pt}
\setlength{\parindent}{0pt}
\setlength{\emergencystretch}{2em}
\newtheorem{theorem}{Theorem}
\newtheorem{corollary}[theorem]{Corollary}
\newtheorem{lemma}[theorem]{Lemma}
\newtheorem{proposition}[theorem]{Proposition}
\theoremstyle{remark}
\newtheorem{remark}[theorem]{Remark}
\DeclareMathOperator{\tr}{tr}
\DeclareMathOperator{\diag}{diag}
\DeclareMathOperator{\rank}{rank}
\newcommand{\E}{\mathbb E}
\newcommand{\R}{\mathbb R}
\newcommand{\C}{\mathbb C}
\newcommand{\F}{\mathrm F}
\newcommand{\sub}[1]{[#1]}
\newcommand{\ind}{\mathbf 1}
\newcommand{\eps}{\epsilon}
\newcommand{\norm}[1]{\left\lVert#1\right\rVert}
\title{Sharp worst-case factors for randomly pivoted Cholesky and LU}
\author{Matthew J. Colbrook\thanks{Department of Applied Mathematics and Theoretical Physics, University of Cambridge, Cambridge, United Kingdom. Email: m.colbrook@damtp.cam.ac.uk.}}
\date{11 September 2026}
\hypersetup{pdfauthor={Matthew J. Colbrook}}
\usepackage{xurl}
\begin{document}
\maketitle

\begin{quote}\small
Authorship is recorded at the submitter's request. This manuscript was developed with AI assistance. Independent agents checked the complete source against RA-02 and RA-03; exact verification was rerun separately. This is not external human peer review or formal certification, and no priority claim is made. The original archive's eligibility hold and failed-access remarks are historical: the live repository checks and submission record accompany this authored edition.
\end{quote}
\begin{abstract}
For every positive integer $r$, an explicit one-parameter family of real,
entrywise-positive, positive-definite matrices of order $r+1$ makes the expected
trace error of $r$ randomly pivoted Cholesky steps approach $2^r$ times the optimal
rank-$r$ trace error. The same family makes the expected squared Frobenius error
of $r$ randomly pivoted LU steps approach $4^r$ times the optimal rank-$r$
squared Frobenius error. Thus both previously established universal factors are
sharp as suprema. The same suprema hold on positive-definite correlation
matrices when the dimension is unrestricted. In particular, a polynomial-in-$r$ factor for same-rank
RPCholesky and a $2^r$ factor for same-rank RPLU are impossible without additional
hypotheses. The proof combines scale-separated Cauchy--Binet expansions with an
exact cancellation over pivot histories. A separate $2\times2$ rational example
already disproves the literal $2^r$ RPLU bound at $r=1$. Exact rational
verification certificates accompany the proof.
\end{abstract}

\section{Algorithms and results}

For a positive-semidefinite matrix $A$, write
\[
 \tau_r(A)=\sum_{j>r}\lambda_j(A),
 \qquad \lambda_1(A)\ge\cdots\ge\lambda_n(A)\ge0.
\]
Randomly pivoted Cholesky starts from $R_0=A$. Given a nonzero residual $R_s$,
it chooses $j$ with probability $(R_s)_{jj}/\tr R_s$ and sets
\begin{equation}\label{eq:chol-update}
 R_{s+1}=R_s-\frac{(R_s)_{:j}(R_s)_{j:}}{(R_s)_{jj}}.
\end{equation}
The residual is positive semidefinite, so its trace is also its nuclear norm.
The known universal bound is
\begin{equation}\label{eq:known-chol}
 \E\tr R_r\le 2^r\tau_r(A)
 \qquad\text{\cite[Lemma 5.5]{CETW}}.
\end{equation}

Randomly pivoted LU starts from $B_0=A$, now allowing a general rectangular
complex matrix. It chooses $(i,j)$ with probability
$|(B_s)_{ij}|^2/\norm{B_s}_{\F}^2$ and sets
\begin{equation}\label{eq:lu-update}
 B_{s+1}=B_s-\frac{(B_s)_{:j}(B_s)_{i:}}{(B_s)_{ij}}.
\end{equation}
A zero entry has zero selection probability. If the residual becomes zero,
the algorithm can stop. If $A_r$ denotes a best rank-$r$ approximation, the
known bound is
\begin{equation}\label{eq:known-lu}
 \E\norm{B_r}_{\F}^2\le4^r\norm{A-A_r}_{\F}^2
 \qquad\text{\cite[Theorem 3]{GW}}.
\end{equation}
Gilles and Wilber conjecture that the factor $4^r$ can be replaced by $2^r$
\cite[discussion after Theorem 3]{GW}. The two questions addressed here are
same-rank approximation questions: exactly $r$ pivots are compared with the
best rank-$r$ error. They correspond to the polynomial-factor RPCholesky
question and the $2^r$ RPLU question catalogued as RA-02 and RA-03,
respectively, in \cite{REPO}. All algorithmic statements below are in exact
arithmetic.

For $r\ge1$, put $n=r+1$. For $0<t<1$, define
\begin{equation}\label{eq:family}
 \begin{split}
 L(t)_{ij}&=
 \begin{cases}
  1,&i=j,\\
  t^j/(i+j),&i>j,\\
  0,&i<j,
 \end{cases}\qquad 1\le i,j\le n,\\
 \eps&=t^{n^2},\qquad
 D=\diag(1,\eps,\ldots,\eps^{n-1}),\qquad
 A_r(t)=L(t)D L(t)^{\mathsf T}.
 \end{split}
\end{equation}
The matrix is positive definite for every $t>0$. Its entries are positive:
for each off-diagonal entry the contribution of the first column of $L$ is
already positive. Rational $t$ gives a rational matrix.

\begin{theorem}[Sharp lower-bound family]\label{thm:main}
Fix $r\ge1$. Apply $r$ steps of the two algorithms to $A=A_r(t)$ in
\eqref{eq:family}. Then
\begin{equation}\label{eq:main-limits}
 \lim_{t\downarrow0}\frac{\E\tr R_r}{\tau_r(A)}=2^r,
 \qquad
 \lim_{t\downarrow0}\frac{\E\norm{B_r}_{\F}^2}
 {\norm{A-A_r^{\mathrm{best}}}_{\F}^2}=4^r,
\end{equation}
where $A_r^{\mathrm{best}}$ denotes a best rank-$r$ approximation of $A$.
Every square submatrix of $A_r(t)$ is nonsingular for all sufficiently small
$t>0$, with the threshold allowed to depend on $r$. Thus every ordered pivot
history appearing in the proof has positive probability.
\end{theorem}

\begin{corollary}[Worst-case constants]\label{cor:sup}
For each $r\ge1$, the supremum of the RPCholesky ratio over all
positive-semidefinite inputs with nonzero rank-$r$ tail is $2^r$. The supremum
of the RPLU squared-error ratio over all matrix inputs with nonzero rank-$r$
error is $4^r$. Restricting the lower-bound examples to real positive-definite,
entrywise-positive matrices does not change either supremum.
\end{corollary}
\begin{proof}
The upper bounds are \eqref{eq:known-chol} and \eqref{eq:known-lu};
Theorem~\ref{thm:main} gives the matching lower bounds.
\end{proof}

\begin{corollary}[Consequences for the proposed improvements]\label{cor:negative}
There are no fixed constants $C>0$ and $p\ge0$ such that
$\E\tr R_r\le Cr^p\tau_r(A)$ for every $r$ and every positive-semidefinite $A$.
For RPLU, no bound of the form
$\E\norm{B_r}_{\F}^2\le C2^r\norm{A-A_r}_{\F}^2$ holds with $C$ independent
of $r$. More generally, no factor $q(r)2^r$, with $q$ a fixed nonnegative polynomial,
can replace $4^r$ uniformly.
\end{corollary}
\begin{proof}
For the first statement, choose $r$ so that $2^{r-1}>Cr^p$, and then use
Theorem~\ref{thm:main} with $t$ sufficiently small. For the second and third
statements, a zero proposed factor is immediately impossible. Otherwise,
$4^r/(2^r q(r))=2^r/q(r)$ is unbounded; again choose $r$ first and then $t$.
No limit uniform in $r$ is needed.
\end{proof}

\section{A two-dimensional RPLU counterexample}

The literal factor-$2^r$ inequality fails even at $r=1$. Take
\begin{equation}\label{eq:small}
 A=\begin{pmatrix}2&1\\1&2\end{pmatrix}.
\end{equation}
Its singular values are $3$ and $1$, so the optimal rank-one squared Frobenius
error is $1$. The total squared entry norm is $10$. Each diagonal pivot has
probability $2/5$ and leaves one nonzero residual entry of magnitude $3/2$.
Each off-diagonal pivot has probability $1/10$ and leaves one nonzero residual
entry of magnitude $3$. Therefore
\begin{equation}\label{eq:small-expectation}
 \E\norm{B_1}_{\F}^2
 =2\left(\frac25\frac94\right)
  +2\left(\frac1{10}\,9\right)
 =\frac{18}{5}>2\norm{A-A_1}_{\F}^2.
\end{equation}
This calculation requires no limiting argument and no numerical eigensolver.

More generally, if every entry of a nonsingular $2\times2$ matrix is nonzero,
a pivot at $(i,j)$ leaves the complementary entry with magnitude
$|\det A|/|A_{ij}|$. Summing over the four pivot probabilities gives
\[
 \E\norm{B_1}_{\F}^2=\frac{4|\det A|^2}{\norm A_{\F}^2},
 \qquad
 \frac{\E\norm{B_1}_{\F}^2}{\sigma_2(A)^2}
 =\frac{4\sigma_1(A)^2}{\sigma_1(A)^2+\sigma_2(A)^2}.
\]
The latter ratio tends to $4$ when the singular-value ratio tends to infinity.
The rest of the note establishes sharpness at every rank, not just at rank one.

\section{Prefix minors and scale separation}

Throughout the proof, $r$ (and hence $n$) is fixed and $t\downarrow0$.
All implicit constants may depend on the fixed dimension. A submatrix indexed
by sets is ordered increasingly unless an appended row or column order is
specified. For $S\subseteq\{1,\ldots,n\}$ with $|S|=s$, set
\[
 M_s(S)=L[S,1\!:\!s],\qquad
 d_s(S)=\sum_{j\in\{1,\ldots,s\}\setminus S}j.
\]
An empty determinant is $1$.

\begin{lemma}[Prefix-row minors]\label{lem:prefix}
For every such $S$,
\begin{equation}\label{eq:prefix-asymptotic}
 \det M_s(S)=\kappa_s(S)t^{d_s(S)}(1+O(t)),
 \qquad \kappa_s(S)\ne0.
\end{equation}
For $s\ge1$, $M_s(S)$ is invertible for sufficiently small $t$, and
$\norm{M_s(S)^{-1}}_2=O(t^{-s})$.
\end{lemma}
\begin{proof}
Let $P=S\cap\{1,\ldots,s\}$, $I=S\setminus\{1,\ldots,s\}$, and
$J=\{1,\ldots,s\}\setminus S$. Thus $|I|=|J|$. Group the rows as $(P,I)$
and the columns as $(P,J)$. Multiply each column indexed by $j\in J$ by
$t^{-j}$, leaving columns in $P$ unchanged. The resulting matrix has a finite
limit of the form
\[
 \begin{pmatrix}I&F\\0&C\end{pmatrix},
 \qquad C=\left[\frac1{i+j}\right]_{i\in I,\,j\in J},
\]
where $F$ is finite. Indeed, $L[P,P]\to I$, $L[I,P]\to0$, and the scaled
$L[I,J]$ equals $C$ exactly. In the block $L[P,J]$, no diagonal entry occurs
because $P$ and $J$ are disjoint, so column scaling leaves bounded entries.

The Cauchy matrix $C$ is nonsingular. For distinct increasing $x_a$ and $y_b$
with positive sums, the elementary Cauchy determinant identity is
\[
 \det\left[\frac1{x_a+y_b}\right]_{a,b=1}^{m}
 =\frac{\prod_{a<a'}(x_{a'}-x_a)\prod_{b<b'}(y_{b'}-y_b)}
 {\prod_{a,b}(x_a+y_b)}\ne0.
\]
The empty block has determinant $1$. Undoing the column scaling gives
\eqref{eq:prefix-asymptotic}, including a possible sign from column ordering.
The scaled matrix and its inverse remain bounded. Undoing the scaling in the
inverse costs at most $t^{-s}$, proving the norm bound.
\end{proof}

\begin{lemma}[Dominance of the leading spectral scales]\label{lem:CB}
For row and column sets $I,J$ of cardinality $s\le n-1$,
\begin{equation}\label{eq:CB-asymptotic}
 \det A[I,J]
 \sim \eps^{s(s-1)/2}\det M_s(I)\det M_s(J).
\end{equation}
For $s=n$, the corresponding equality is exact, with $M_n=L$ and
$\det L=1$. Consequently, every square minor of $A$ is nonzero for all
sufficiently small positive $t$.
\end{lemma}
\begin{proof}
Cauchy--Binet gives
\[
 \det A[I,J]=\sum_{\substack{T\subseteq\{1,\ldots,n\}\\|T|=s}}
 \det L[I,T]\det L[J,T]\,
 \eps^{\sum_{j\in T}(j-1)}.
\]
The smallest power of $\eps$ occurs uniquely at $T=\{1,\ldots,s\}$ and is
$s(s-1)/2$. All other powers are at least one larger. Since $L(t)$ remains
bounded, all minors in the finite sum remain bounded. Thus
\[
 \det A[I,J]=\eps^{s(s-1)/2}
 \left(\det M_s(I)\det M_s(J)+O(\eps)\right).
\]
Lemma~\ref{lem:prefix} shows that the leading product has nonzero coefficient
and $t$-degree at most $s(s+1)$. Since $\eps=t^{n^2}$ and
$n^2>s(s+1)$ for $s\le n-1$, the remainder divided by that product tends to
zero. If $s=n$, the sum has only one term. There are finitely many minors in
a fixed dimension, so their eventual nonvanishing holds simultaneously.
\end{proof}

\section{Residual normalizers}

For $|S|=s<n$, let $l_{s+1}=L[:,s+1]$ and define
\begin{equation}\label{eq:w}
 c_s(S)=M_s(S)^{-1}l_{s+1}[S],\qquad
 w_s(S)=l_{s+1}-L[:,1\!:\!s]c_s(S),\qquad
 h_s(S)=\norm{w_s(S)}_2^2.
\end{equation}
For $s=0$, the coefficient vector is empty and $w_0(\varnothing)=l_1$.
The vector $w_s(S)$ vanishes on $S$.

\begin{lemma}[A binary limiting rule]\label{lem:h}
For each fixed $S$ of size $s<n$,
\begin{equation}\label{eq:h-indicator}
 \lim_{t\downarrow0}\frac1{h_s(S)}
 =\ind_{\{s+1\notin S\}}.
\end{equation}
Moreover, $h_s(S)\ge1/4$ for all sufficiently small $t$.
\end{lemma}
\begin{proof}
Because $L\to I$, its smallest singular value is at least $1/2$ eventually.
The vector $w_s(S)$ equals $L$ times a vector whose first $s$ entries are
$-c_s(S)$, whose $(s+1)$st entry is $1$, and whose remaining entries vanish.
It follows that
\begin{equation}\label{eq:h-lower}
 h_s(S)\ge\frac14\left(1+\norm{c_s(S)}_2^2\right).
\end{equation}
If $s+1\notin S$, every entry of $l_{s+1}[S]$ is either zero or
$O(t^{s+1})$. Lemma~\ref{lem:prefix} therefore gives $c_s(S)=O(t)$, so
$w_s(S)\to e_{s+1}$ and $h_s(S)\to1$.

If $s+1\in S$, the row of $M_s(S)$ indexed by $s+1$ has norm $O(t)$, whereas
the same entry of $l_{s+1}[S]$ is exactly $1$. The equation defining $c_s(S)$
forces $\norm{c_s(S)}_2\ge c/t$ for a positive constant $c$ and small $t$.
Equation~\eqref{eq:h-lower} then gives $h_s(S)\to\infty$.
The empty-set case is included in the first alternative.
\end{proof}

For equally sized row and column sets $I,J$, define the full, zero-padded
interpolatory residual
\[
 \mathcal R(I,J)=A-A[:,J]A[I,J]^{-1}A[I,:].
\]
For empty sets this means $\mathcal R(\varnothing,\varnothing)=A$.
Its rows in $I$ and columns in $J$ vanish. After any nonzero-pivot history with
selected row set $I$ and column set $J$, Gaussian elimination gives this
residual, independently of the order within the selected sets. In the
Cholesky case $I=J$.

\begin{lemma}[Asymptotic normalizers]\label{lem:normalizers}
For $|I|=|J|=s<n$, one has
\begin{align}
 T_s(S):=\tr\mathcal R(S,S)&\sim\eps^s h_s(S),\label{eq:trace-asymptotic}\\
 Q_s(I,J):=\norm{\mathcal R(I,J)}_{\F}^2
 &\sim\eps^{2s}h_s(I)h_s(J).\label{eq:norm-asymptotic}
\end{align}
In particular,
\begin{align}
 \frac{\eps^s}{T_s(S)}&\longrightarrow\ind_{\{s+1\notin S\}},\label{eq:trace-indicator}\\
 \frac{\eps^{2s}}{Q_s(I,J)}
 &\longrightarrow\ind_{\{s+1\notin I\}}\ind_{\{s+1\notin J\}}.
 \label{eq:norm-indicator}
\end{align}
\end{lemma}
\begin{proof}
For $i\notin I$ and $j\notin J$, the bordered determinant identity, with $i$
and $j$ appended to the existing orders, gives
\[
 \mathcal R(I,J)_{ij}
 =\frac{\det A[(I,i),(J,j)]}{\det A[I,J]}.
\]
Apply Lemma~\ref{lem:CB} to numerator and denominator. The ratio of the
powers of $\eps$ is $\eps^s$. A second bordered determinant identity gives
\[
 \det L[(I,i),1\!:\!s+1]=\det M_s(I)\,w_s(I)_i,
\]
and the analogous column identity gives $\det M_s(J)\,w_s(J)_j$.
Consequently,
\begin{equation}\label{eq:entry-asymptotic}
 \mathcal R(I,J)_{ij}
 =\eps^s w_s(I)_i w_s(J)_j(1+o(1)).
\end{equation}
The two components in this formula are nonzero for sufficiently small $t$,
by Lemma~\ref{lem:prefix} applied to the bordered prefix minors. For a fixed
dimension, the relative $o(1)$ is uniform over the finitely many relevant
entries and sets.

For principal residuals, the diagonal leading terms in
\eqref{eq:entry-asymptotic} are $\eps^s w_s(S)_i^2$. Summing them proves
\eqref{eq:trace-asymptotic}. Squaring and summing all entries proves
\eqref{eq:norm-asymptotic}. These are relative asymptotics even when a vector
$w_s$ diverges: the maximum relative error tends to zero before the weighted
finite sum is taken. Finally use Lemma~\ref{lem:h}; its lower bound on $h_s$
justifies inversion also in the diverging cases.
\end{proof}

\section{Exact cancellation over pivot histories}

The next identities hold before taking any limit. Recall that $n=r+1$.
For a Cholesky history $(j_1,\ldots,j_r)$ of distinct indices, let
$S_s=\{j_1,\ldots,j_s\}$. For an LU history, define similarly
$I_s=\{i_1,\ldots,i_s\}$ and $J_s=\{j_1,\ldots,j_s\}$.

\begin{lemma}[History identities]\label{lem:paths}
For an $(r+1)\times(r+1)$ positive-definite input with all square minors
nonzero, the expectations after $r$ pivots satisfy
\begin{align}
 \E\tr R_r
 &=\det A\sum_{(j_1,\ldots,j_r)}
     \prod_{s=0}^{r-1}\frac1{T_s(S_s)},\label{eq:chol-path}\\
 \E\norm{B_r}_{\F}^2
 &=|\det A|^2\sum_{\substack{(i_1,\ldots,i_r)\\(j_1,\ldots,j_r)}}
     \prod_{s=0}^{r-1}\frac1{Q_s(I_s,J_s)}.\label{eq:lu-path}
\end{align}
The sums range over ordered lists of distinct indices; in the second sum
row indices and column indices are separately distinct.
\end{lemma}
\begin{proof}
In Cholesky, the product of the $r$ pivot values is the principal determinant
$\det A[S_r,S_r]$. The probability of the history is this product divided by
$\prod_{s=0}^{r-1}T_s(S_s)$. Only one diagonal residual entry remains after
$r$ pivots, and its value is $\det A/\det A[S_r,S_r]$. The determinant factors
cancel when probability is multiplied by final trace.

In LU, the product of the pivot values equals the determinant of the selected
$r\times r$ block in the ordered row and column lists. Its squared modulus
is the numerator of the history probability. The only remaining residual
entry has squared modulus $|\det A|^2/|\det A[I_r,J_r]|^2$; permutations
change only a sign before taking the modulus. Cancellation gives
\eqref{eq:lu-path}. Summing history contributions proves both identities.
No independence of the LU row and column choices is being assumed.
\end{proof}

\section{Counting the surviving histories}

First identify the optimal error scale. The determinant and inverse of the
family satisfy
\begin{equation}\label{eq:det-inverse}
 \det A=\eps^{r(r+1)/2},\qquad
 \eps^r A^{-1}
 =L^{-\mathsf T}\diag(\eps^r,\eps^{r-1},\ldots,1)L^{-1}
 \longrightarrow e_n e_n^{\mathsf T}.
\end{equation}
The largest eigenvalue on the right-hand side tends to $1$, whence
\begin{equation}\label{eq:lambda-scale}
 \frac{\lambda_n(A)}{\eps^r}\longrightarrow1.
\end{equation}
Since $A$ has order $r+1$ and is positive definite,
$\tau_r(A)=\lambda_n(A)$ and
$\norm{A-A_r^{\mathrm{best}}}_{\F}^2=\lambda_n(A)^2$.

Using \eqref{eq:det-inverse} in the first history identity gives
\begin{equation}\label{eq:normalized-chol-path}
 \frac{\E\tr R_r}{\lambda_n(A)}
 =\frac{\eps^r}{\lambda_n(A)}
 \sum_{(j_1,\ldots,j_r)}\prod_{s=0}^{r-1}
       \frac{\eps^s}{T_s(S_s)}.
\end{equation}
By Lemma~\ref{lem:normalizers}, the product for a fixed history tends to $1$
precisely when
\begin{equation}\label{eq:survive}
 s+1\notin S_s\qquad(0\le s\le r-1),
\end{equation}
and otherwise tends to $0$. The condition at $s=0$ is vacuous.

Condition~\eqref{eq:survive} is equivalent to
\begin{equation}\label{eq:twochoices}
 j_u\in\{1,\ldots,u,n\}\setminus\{j_1,\ldots,j_{u-1}\}
 \qquad(1\le u\le r).
\end{equation}
Indeed, choosing a label $q\in\{2,\ldots,r\}$ before position $q$ violates
\eqref{eq:survive} at $s=q-1$. Conversely, if every such label is chosen no
earlier than its own position, no prefix can contain its next label. Label
$n=r+1$ has no corresponding normalizer constraint, because the normalizers
stop at $s=r-1$.

At position $u$, the allowed set before removing earlier choices has $u+1$
elements. All $u-1$ earlier choices lie in it. Thus there are exactly two
choices at every position, and exactly $2^r$ surviving histories.
The sum in \eqref{eq:normalized-chol-path} is finite for each fixed $r$.
Together with \eqref{eq:lambda-scale}, this proves the Cholesky limit.

For LU, the corresponding normalized identity is
\begin{equation}\label{eq:normalized-lu-path}
 \frac{\E\norm{B_r}_{\F}^2}{\lambda_n(A)^2}
 =\frac{\eps^{2r}}{\lambda_n(A)^2}
 \sum_{\substack{(i_1,\ldots,i_r)\\(j_1,\ldots,j_r)}}
 \prod_{s=0}^{r-1}\frac{\eps^{2s}}{Q_s(I_s,J_s)}.
\end{equation}
Lemma~\ref{lem:normalizers} says that a summand tends to $1$ exactly when
both the row history and the column history satisfy \eqref{eq:survive}; it
otherwise tends to $0$. Each list has $2^r$ possibilities, giving $4^r$
ordered pairs. This is a factorization of the limiting indicator and a
combinatorial count, not a claim that the algorithm samples its row and column
histories independently. Equation~\eqref{eq:lambda-scale} now proves the LU
limit and completes the proof of Theorem~\ref{thm:main}.
\qed

\section{Unit-diagonal inputs by row replication}

The examples of order $r+1$ are not unit-diagonal. Nevertheless, requiring a
unit diagonal does not improve either worst-case constant when arbitrary
matrix dimensions are allowed.

\begin{proposition}[Correlation-matrix sharpness]\label{prop:correlation}
For each fixed $r\ge1$, both suprema in Corollary~\ref{cor:sup} remain unchanged
when inputs are restricted to real, entrywise-positive, positive-definite
matrices with every diagonal entry equal to $1$.
\end{proposition}
\begin{proof}
First let $A$ be any rational positive-definite matrix of order $n$ with
positive entries. Choose $c>0$ rational so that
$m_i=A_{ii}/c$ is a positive integer for every $i$. Put $N=\sum_i m_i$.
Partition the $N$ row labels into groups of sizes $m_1,\ldots,m_n$, and
construct $P\in\R^{N\times n}$ by assigning each row in group $i$ the vector
$e_i^{\mathsf T}/\sqrt{m_i}$. Then $P^{\mathsf T}P=I_n$, and
\[
 C=c^{-1}PAP^{\mathsf T}
\]
is positive semidefinite, entrywise positive, and unit-diagonal. Its nonzero
eigenvalues are those of $A/c$.

The two pivoting processes on $C$ exactly reproduce those on $A$ after
aggregating pivot labels by group. More precisely, if the current residual
has the form $c^{-1}PBP^{\mathsf T}$, a Cholesky pivot in group $i$ has
aggregate probability
\[
 \frac{m_i B_{ii}/(cm_i)}{\tr B/c}=\frac{B_{ii}}{\tr B}.
\]
For LU, the aggregate probability of a row in group $i$ and a column in
group $j$ is
\[
 \frac{m_i m_j |B_{ij}|^2/(c^2m_i m_j)}{\norm B_{\F}^2/c^2}
 =\frac{|B_{ij}|^2}{\norm B_{\F}^2}.
\]
Substituting either pivot into its rank-one update shows that the new
residual is $c^{-1}P B_{\mathrm{new}}P^{\mathsf T}$. These assertions also
respect zero pivot probabilities. Thus the trace error is scaled by $1/c$
and the squared Frobenius error by $1/c^2$, exactly as are the respective
optimal rank-$r$ errors. Both expected-error ratios are unchanged.

To obtain positive definiteness, put
\[
 C_\delta=\frac{C+\delta I_N}{1+\delta},\qquad \delta>0.
\]
This matrix is positive definite, entrywise positive, and unit-diagonal.
Fix a finite stopping rank $r$ for which the optimal error of $C$ is positive.
Every pivot history having positive probability for $C$ has nonzero pivots
at all of its steps. Along each such history, the pivot values, probabilities,
and terminal residuals for $C_\delta$ converge to those for $C$, because the
rank-one updates are rational functions whose denominators remain nonzero.
There are only finitely many histories. Summing these convergent contributions
and discarding all other, nonnegative contributions proves
\[
 \liminf_{\delta\downarrow0} F_r(C_\delta)\ge F_r(C),
\]
where $F_r$ is either expected-error ratio. Here the optimal-error denominators
converge to their positive values for $C$, by continuity of the eigenvalues.
Full continuity of the LU expectation at zero pivot entries is not needed.

Finally choose rational $t$ in Theorem~\ref{thm:main} to approach the desired
sharp constant, apply the replication construction to $A_r(t)$, and then
take $\delta$ sufficiently small. The known universal upper bounds complete
the proof. The replicated dimension $N$ is not bounded uniformly in $t$.
\end{proof}

\section{Exact finite certificates and scope}

The limiting proof establishes a counterexample for every proposed smaller
uniform factor. The accompanying code also supplies explicit finite rational
certificates. These are useful because the matrices are deliberately
ill-conditioned, making an ordinary double-precision eigensolve unreliable.

Let $v=L^{-\mathsf T}e_n$ and $\nu=\norm v_2^2$. From
\eqref{eq:det-inverse},
\[
 \eps^r A^{-1}\succeq vv^{\mathsf T},\qquad
 \lambda_n(A)\le\frac{\eps^r}{\nu}.
\]
For $0\le s\le r-1$, let $\mathcal G_s$ consist of all $s$-element subsets
of $\{1,\ldots,s,n\}$, and define the exact rational maxima
\[
 m_s=\max_{S\in\mathcal G_s}\frac{T_s(S)}{\eps^s},\qquad
 q_s=\max_{I,J\in\mathcal G_s}\frac{Q_s(I,J)}{\eps^{2s}}.
\]
Every history counted in \eqref{eq:twochoices} has its $s$th prefix in
$\mathcal G_s$. Retaining only these histories in
\eqref{eq:normalized-chol-path} and \eqref{eq:normalized-lu-path} gives
\begin{equation}\label{eq:cert}
 \frac{\E\tr R_r}{\lambda_n(A)}\ge
 \frac{2^r\nu}{\prod_{s=0}^{r-1}m_s},\qquad
 \frac{\E\norm{B_r}_{\F}^2}{\lambda_n(A)^2}\ge
 \frac{4^r\nu^2}{\prod_{s=0}^{r-1}q_s}.
\end{equation}
For a finite certificate it suffices to verify that these retained pivot
histories are valid; nonsingularity of their prefix and terminal blocks is
checked exactly. Omitted histories contribute nonnegative error.

The file \texttt{exact\_certificate.py} evaluates \eqref{eq:cert} with Python
\texttt{Fraction} arithmetic. The separate file
\texttt{exact\_enumeration.py} recursively executes all actual rank-one pivot
updates for small dimensions, without using the history cancellation formula
to compute expectations. It encloses the optimal error using
\[
 \frac{\eps^r}{\tr(\eps^r A^{-1})}\le\lambda_n(A)
 \le\frac{\eps^r}{\nu}.
\]
Both endpoints and both error expectations are rational, so the resulting
ratio enclosures are rigorous. Selected lower bounds at $t=1/100$ are shown
below; the displayed decimals are rounded downward, and the exact fractions
are stored in the accompanying JSON records.

\begin{center}
\small
\begin{tabular}{@{}r l r r r r@{}}
\toprule
$r$ & Certificate & Cholesky ratio $>$ & $2^r$ & LU ratio $>$ & $4^r$\\
\midrule
2 & All pivot updates & 3.99995644 & 4 & 15.99965154 & 16\\
3 & All pivot updates & 7.99957072 & 8 & 63.99313178 & 64\\
8 & Retained histories & 255.85420697 & 256 & 65461.37522779 & 65536\\
\bottomrule
\end{tabular}
\end{center}

The rank-eight certificate uses $n=9$ and $\eps=(1/100)^{81}$.
The unusually small scales are part of an exact-arithmetic worst-case
construction, not a proposed practical test in double precision.

The small-dimensional family has unrestricted diagonal entries;
Proposition~\ref{prop:correlation} separately covers the unit-diagonal setting
in unrestricted dimension. No bounded-condition-number claim is made.
These results do not resolve how many pivots beyond $r$ suffice for a
$(1+\delta)$-relative trace approximation, the separate oversampling question
RA-01. In particular, after $r+1$ pivots the matrices in this note are recovered
exactly. Zero-padding embeds the lower-bound family in any larger dimension
without altering its nonzero-pivot process, but does not turn this same-rank
result into an oversampling lower bound.

\begin{thebibliography}{9}
\bibitem{CETW}
Y.~Chen, E.~N. Epperly, J.~A. Tropp, and R.~J. Webber.
\emph{Randomly pivoted Cholesky: Practical approximation of a kernel matrix
with few entry evaluations}.
Communications on Pure and Applied Mathematics, 78(5):995--1041, 2025.
\href{https://doi.org/10.1002/cpa.22234}{doi:10.1002/cpa.22234}.
The universal same-rank bound is Lemma~5.5.

\bibitem{GW}
M.~A. Gilles and H.~Wilber.
\emph{Low-Rank Approximation by Randomly Pivoted LU}.
arXiv:2601.22344, 2026.
\url{https://arxiv.org/abs/2601.22344}.
The equation numbering here follows the
\href{https://www.researchgate.net/publication/400339985_Low-Rank_Approximation_by_Randomly_Pivoted_LU}{author-posted full text}
dated 2 February 2026. The sampling rule is equation~(3), the bound is
Theorem~3, and the proposed $2^r$ improvement appears immediately after it.

\bibitem{REPO}
\emph{Open Problems in Numerical Linear Algebra}, randomized and low-rank
approximation category, entries RA-02 and RA-03.
\url{https://github.com/ajt60gaibb/OpenProblemsInNLA/tree/main/randomized-and-low-rank-approximation}.
The repository-audit record accompanying this note distinguishes the retrieved
category listing from source-file and pull-request checks that could not be
completed. The mathematical statements proved here are specified in full
above and do not depend on a repository status label.
\end{thebibliography}
\end{document}
